\section{Le téléphone (1876--1940)}\label{sec:telephone}
% =============================================================================


% §1.5 — DRAFT v1 — 26 avril 2026
% Structure en quatre temps :
%   (1) Ouverture : le non-codifiable et le déplacement du canal
%   (2) Bell, Gray et le brevet : focusing device inversé
%   (3) Le monopole et les modèles institutionnels (Vail, 4 pays)
%   (4) La commutation et les demoiselles (Feigenbaum & Gross)
% Bilan : tableau D1–D10 + 0-TER
% Sources : VdK Telephone, Bertho 1984 (pp. 51–79), Feigenbaum & Gross
%   (QJE 2024, Management Science 2024), Mueller 1993 (via VdK)

Au milieu des années~1870, l'industrie du télégraphe est la plus puissante des industries de communication. Western Union, constituée en~1856, est l'une des premières grandes entreprises américaines~; elle possède des dizaines de milliers de kilomètres de lignes, emploie des milliers d'opérateurs, et génère des profits considérables. En Europe, les administrations d'État (Prusse, France, Angleterre) contrôlent des réseaux tout aussi étendus. C'est dans ce terreau industriel~--- dans ses ateliers, avec ses ingénieurs, contre ses problèmes~--- que naît le téléphone.

Le problème qui mobilise les meilleurs esprits de cette industrie est celui de la saturation des lignes. Les conducteurs de cuivre coûtent cher, et chaque fil ne porte qu'un message à la fois. Faire passer plusieurs messages simultanément sur un même fil~--- le multiplexage~--- est le graal technique des années~1870. Edison y travaille avec le quadruplex (1874), qui permet quatre transmissions simultanées. Elisha Gray, ingénieur chez Western Union depuis~1870, y travaille avec la télégraphie harmonique, qui exploite la superposition de fréquences. Ce sont des hommes du télégraphe, formés par le télégraphe, qui raisonnent dans les catégories du télégraphe~: la vitesse, le débit, le nombre de messages par fil.

Mais l'idée de transmettre la voix humaine par un circuit électrique est plus ancienne que cette course, et ne vient pas seulement du monde anglo-saxon. En~1854, Charles Bourseul, employé de l'administration française des télégraphes, avait décrit dans un article de \emph{L'Illustration} le principe d'un appareil capable de transmettre la parole, sans jamais parvenir à le réaliser\footnote{Bourseul écrivait~: \og Supposez que l'on parle près d'une plaque mobile assez flexible pour ne perdre aucune des vibrations produites par la voix, que cette plaque établisse et interrompe successivement la communication avec une pile~: vous pourrez avoir à distance une autre plaque qui exécutera exactement les mêmes vibrations.\fg{} Cité par \textcite{bertho_histoire_1984}, p.~51.}. En~1861, l'Allemand Philipp Reis construisit un appareil qu'il appela \emph{Telephon}, capable de transmettre des sons musicaux et, dans des conditions favorables, des fragments de parole~; il le présenta devant la Société de physique de Francfort, mais ne parvint pas à intéresser les industriels du télégraphe, qui n'y virent qu'un jouet acoustique\footnote{Reis mourut en~1874, deux ans avant le brevet de Bell, sans avoir obtenu de reconnaissance pour son travail. L'appareil de Reis transmettait le son par variation de résistance d'un contact, pas par variation continue de courant comme le ferait celui de Bell~; la distinction technique est réelle mais n'enlève rien à l'antériorité du concept.}. Antonio Meucci, émigré italien installé à New York, déposa un \emph{caveat} en~1871 pour un dispositif de communication vocale, mais ne put payer les frais de renouvellement et perdit sa protection en~1874.

\dimmark{D10}{+}%
Ces échecs ne sont pas des accidents biographiques. Ce qu'ils révèlent est un fait structurel~: le concept de transmission de la voix existait depuis au moins vingt ans, dans trois pays et trois langues, mais il ne trouvait pas de véhicule institutionnel capable de le porter. L'industrie du télégraphe, seule infrastructure de communication à longue distance, n'en voulait pas~; les inventeurs isolés (Bourseul, Reis, Meucci) n'avaient ni les capitaux ni les réseaux d'influence nécessaires pour transformer un prototype en système commercial. Il fallut que deux hommes issus de cette industrie~--- mais chacun à sa manière en marge d'elle~--- arrivent au même point presque au même moment.


\subsection{Bell, Gray et le brevet~174,465}\label{subsec:bell}

Le 14~février~1876, deux demandes concurrentes arrivèrent au bureau des brevets de New York à quelques heures d'intervalle. Alexander Graham Bell, par l'entremise de son avocat Anthony Pollok, déposa une demande de brevet pour un appareil de transmission de la voix par courant électrique. Elisha Gray, par l'entremise de son avocat William Baldwin, déposa un \emph{caveat}~--- une déclaration d'intention de breveter~--- pour un dispositif de même nature\footnote{La chronologie exacte a fait l'objet de controverses durables. \textcite{evenson_telephone_2000} a montré que le \emph{caveat} de Gray, bien qu'arrivé physiquement au bureau des brevets avant la demande de Bell, fut enregistré dans le journal quotidien (\emph{cash blotter}) après celle-ci, en raison d'une procédure accélérée demandée par Pollok. La question de savoir si Bell avait eu connaissance du contenu du \emph{caveat} de Gray par l'intermédiaire de l'examinateur Zenas Fisk Wilber reste débattue~; voir \textcite{shulman_telephone_2008} pour l'enquête la plus récente.}. Le brevet de Bell, n\textsuperscript{o}~174,465, fut accordé le 7~mars~1876.

Gray était un professionnel du télégraphe, en poste chez Western Union depuis~1870, rompu aux problèmes du multiplexage. Bell était un outsider~: issu d'une famille d'émigrés anglais installée au Canada, spécialisée depuis trois générations dans l'éducation des sourds, il avait étudié les problèmes de l'audition et de la voix humaine avant de s'intéresser à l'électricité \parencite{bertho_histoire_1984}. S'il cherchait un télégraphe harmonique, c'est qu'il savait que les compagnies télégraphiques feraient la fortune de son inventeur~; mais il devait lutter contre ses propres commanditaires, Gardiner Hubbard et Thomas Sanders, pour se détourner du multiplexage et se consacrer au téléphone, dont aucun prototype ne fonctionnait encore au moment du dépôt du brevet. La tension entre ce que le financeur demande (le multiplexage, dont la valeur est immédiate et mesurable) et ce que l'inventeur poursuit (la voix, dont la valeur est incertaine) est une instance de la bidirectionnalité notée dans le préambule méthodologique~: la demande oriente la technique, mais la technique la plus transformatrice peut être celle qui ne répond à aucune demande formulée.

\dimmark{D5}{?}%
Ce qui importe pour l'analyse n'est pas la question de priorité, que les tribunaux tranchèrent en faveur de Bell après six cents procès étalés sur dix-huit ans, mais deux faits que l'épisode met en lumière. Le premier est un fait juridique. En~1881, dans l'affaire \emph{Spencer}, le juge Lowell statua que Bell avait créé un \og \emph{new art}\fg{}~--- une technique entièrement nouvelle~--- et qu'il avait droit à la revendication la plus large possible sur la transmission électrique de la parole\footnote{\og Bell had discovered a new art~--- that of transmitting speech by electricity~--- and has a right to hold the broadest claim for it which can be permitted in any case\fg{} (\textcite{beauchamp_who_2010}, p.~863).}. Cette qualification de \og nouvel art\fg{} transformait un brevet d'appareil en brevet de principe, couvrant non pas un dispositif particulier mais la totalité du champ technique. C'est ce brevet de principe qui, pendant dix-huit ans, donna à Bell et à ses successeurs un monopole légal sur la téléphonie aux États-Unis. Le mécanisme de monopole qui en résultait~--- une barrière juridique fondée sur la portée du brevet, non sur les rendements de réseau ni sur le verrouillage d'un consommable~--- était distinct de ceux qu'avaient fait apparaître le télégraphe et la tabulatrice.

\dimmark{D5}{?}%
Le second fait est un fait de perception. Gray, qui travaillait pour Western Union depuis~1870 et avait résolu plusieurs problèmes du multiplexage télégraphique, ne voyait dans le téléphone qu'une curiosité~: \og Le téléphone parlant, écrivait-il le lendemain à son associé, est un bel objet du point de vue scientifique\ldots{} mais si on le considère sous l'angle commercial, il n'a aucune valeur. Avec un fil on peut faire actuellement beaucoup plus, et la vitesse est la seule chose qui nous intéresse\fg{}\footnote{Cité par \textcite{bertho_histoire_1984}, p.~53. La formulation anglaise, rapportée par \textcite{hounshell_elisha_1975}, est~: \og \ldots talking telegraph is a beautiful thing in a scientific point of view\ldots{} But if you look at it in a business light it is of no importance.\fg{}}. La même année, William Orton, président de Western Union, refusa d'acheter les brevets de Bell pour cent mille dollars\footnote{\textcite{hounshell_elisha_1975}, p.~157. Orton voulait un télégraphe multiplex, pas un \og scientific curiosity\fg{}. Western Union changea d'avis l'année suivante et tenta de contester les brevets de Bell, sans succès.}. Ce que les professionnels du télégraphe ne parvenaient pas à voir, c'est que la valeur du téléphone résidait précisément dans ce que le télégraphe ne pouvait pas faire~: transmettre le non-codifiable. Leur grille d'évaluation~--- le nombre de messages par fil, la vitesse de transmission, le coût par mot~--- était celle d'un système fondé sur le codage, et elle rendait invisible la valeur d'un système fondé sur la voix. C'est un cas de ce que \textcite{rosenberg_perspectives_1976} appelle un \emph{focusing device} inversé~: le paradigme technique en place concentre l'attention sur certains problèmes (le multiplexage, la vitesse) au point d'aveugler les praticiens sur des possibilités qui relèvent d'un autre registre.

\dimmark{D6/D7}{+}%
Ce que le téléphone rend transmissible, c'est précisément ce que le codage télégraphique détruisait~: le ton d'une négociation, la crédibilité d'un interlocuteur, l'urgence d'une instruction. Mais la portée de cette rupture ne se manifesta pas immédiatement. \textcite{ploeckl_market_2023}, étudiant la diffusion du téléphone en Bavière entre~1883 et~1905 à travers 306~centraux téléphoniques, a montré que le téléphone était d'abord une commodité urbaine locale~: la taille de la population de la ville~--- et non les connexions avec d'autres villes~--- dominait largement la décision d'ouvrir un central et le nombre d'abonnés. Les premiers utilisateurs s'appelaient dans la même ville, pas à longue distance~; la présence d'une gare n'accélérait pas la diffusion et avait même un effet négatif sur l'adoption, ce qui signale que le téléphone, contrairement au télégraphe, naît hors de toute filiation avec le réseau ferroviaire. Ce n'est que dans les villes plus petites et plus tardives que l'accès aux réseaux extérieurs commença à peser sur l'adoption. Le téléphone ne déplaça donc pas d'emblée le canal par lequel la transmission agit sur l'économie~; il commença par quelque chose que le télégraphe ne faisait pas du tout~: la communication de proximité, intra-urbaine, quotidienne. Le déplacement vers la coordination organisationnelle à longue distance~--- l'AT\&T de Vail, la ligne transcontinentale de~1915~--- fut un processus de trente ans, pas un fait contemporain de l'invention\footnote{\textcite{juhasz_spinning_2018} ont proposé la notion de filtre de codifiabilité pour caractériser ce que les techniques de transmission rendent échangeable à distance. \textcite{steinwender_real_2018} a montré que le télégraphe déplaçait cette frontière en créant un marché unifié du coton entre Liverpool et New York. Le téléphone la déplace une seconde fois, mais la conséquence économique de ce déplacement~--- l'action sur les organisations plutôt que sur les marchés~--- ne se manifeste qu'une fois le réseau longue distance construit. Cette séquence sera reprise dans le bilan de cette section.}. Et c'est précisément ce décalage temporel qui éclaire le jugement de Gray~: en~1876, le téléphone était effectivement sans valeur dans les catégories du télégraphe, parce que sa valeur résidait dans un usage~--- la communication locale~--- que ces catégories ne savaient pas mesurer.


\subsection{Le Bell System et la construction du monopole}\label{subsec:bell_system}

Le brevet de Bell expirait en~1894. Les dix-huit années qui le séparaient de son dépôt furent employées par ses successeurs à construire, autour de la protection légale, un système intégré dont la logique dépassait de loin celle du brevet lui-même. La \emph{Bell Telephone Company}, créée en juillet~1877 par Bell et ses associés Gardiner Hubbard et Thomas Sanders, était une petite entreprise dont les ressources ne permettaient ni de fabriquer les appareils en série ni d'exploiter les réseaux. Le modèle choisi fut celui de la concession~: Bell accordait à des opérateurs locaux le droit d'exploiter le service téléphonique sur un territoire donné, en échange d'une part de leur capital et de l'obligation de louer (non d'acheter) les appareils Bell. En~1880, les trois cents opérateurs locaux ainsi créés desservaient environ 48\,000~abonnés, concentrés dans les grandes villes\footnote{Données de \textcite{joel_history_1975}, reprises par VdK. En~1895, 57\,\% des téléphones étaient dans des villes de plus de 50\,000~habitants, et 3\,\% seulement dans les zones rurales \parencite{mueller_universal_1993}.}.

\dimmark{D2}{*}%
En~1879, un accord mit fin à la première bataille entre Bell et Western Union~: Bell obtenait le marché du téléphone, Western Union conservait celui du télégraphe. Cet accord, négocié après que les tribunaux eurent reconnu la validité des brevets Bell, eut pour effet de séparer les deux industries de communication, qui partageaient pourtant les mêmes poteaux, les mêmes fils de cuivre et, dans une large mesure, les mêmes ingénieurs. La séparation n'était pas technique~; elle était juridico-commerciale, et elle allait durer un siècle.

C'est Theodore Vail qui donna au système sa structure durable. Vail n'était pas un technicien~: il venait du service postal américain, où il avait réorganisé l'acheminement du courrier par voie ferrée \parencite{bertho_histoire_1984}. C'était, selon la formule de Bertho, \og un stratège qui pense en termes de réseaux\fg{}. Sous sa direction, trois décisions architecturèrent le monopole. La première fut de déplacer la compagnie à New York et de lui donner des statuts permettant la prise de participations dans les compagnies locales~: les opérateurs cessèrent d'être des licenciés indépendants pour devenir des filiales dont la maison mère définissait la politique d'ensemble. La seconde fut de racheter, en~1882, la \emph{Western Electric Manufacturing Company}~--- alors la plus grande usine de matériel électrique des États-Unis, fondée par Elisha Gray et Enos Barton~--- pour intégrer la production du matériel. Aucun autre pays ne connaîtrait cette intégration de l'opérateur et du fabricant dans une même structure \parencite{bertho_histoire_1984}. La troisième fut de créer, en~1885, l'\emph{American Telephone and Telegraph Company} (AT\&T) pour les liaisons longue distance entre les réseaux locaux~: celui qui contrôlait l'interurbain contrôlait l'accès de chaque réseau local au reste du monde.

\dimmark{D2}{*}%
Le mécanisme de monopole qui en résultait procédait d'une logique distincte de celles observées dans les cas précédents. Le télégraphe avait convergé vers le monopole par les rendements de réseau~: chaque nouvel abonné augmentait la valeur du réseau pour tous les autres, et le premier opérateur à atteindre la masse critique éliminait ses concurrents (§\,\ref{sec:telegraph}). La tabulatrice avait convergé vers le monopole par le verrouillage du consommable~: les cartes perforées liaient le client à IBM indépendamment de la qualité de la machine (§\,\ref{sec:mechanography}). Le téléphone convergeait vers le monopole par l'intégration verticale et le contrôle de l'interconnexion~: Vail ne se contentait pas de posséder le réseau le plus étendu~; il possédait aussi l'usine qui fabriquait les appareils, le réseau longue distance qui connectait les villes entre elles, et le portefeuille de brevets qui interdisait aux concurrents d'utiliser les mêmes composants. Six cents procès en contrefaçon maintinrent cette barrière pendant dix-huit ans.

\emergence{D2~--- Le monopole du téléphone procède d'un troisième mécanisme, distinct des rendements de réseau (télégraphe) et du verrouillage du consommable (tabulatrice)~: l'intégration verticale (fabricant + opérateur + longue distance) combinée au contrôle de l'interconnexion. Trois technologies de communication, trois mécanismes de monopole, un même résultat~--- la concentration du marché. \questionch{2}{D2~--- Si le résultat (monopole) est stable mais le mécanisme varie, qu'est-ce qui détermine le choix du mécanisme~? La nature de la technologie, la structure institutionnelle, ou le hasard historique~?}}


\subsection{Quatre pays, quatre trajectoires, un résultat}\label{subsec:quatre_pays}

L'expérience américaine n'est qu'un des quatre modèles institutionnels par lesquels le téléphone fut organisé dans les pays industrialisés. Leur comparaison révèle que la convergence vers l'opérateur unique n'est pas un fait technique mais un fait politique dont la forme varie selon le contexte.

\dimmark{D2/D7}{*}%
En Prusse, le téléphone fut immédiatement placé sous l'autorité de l'administration des postes, sans phase de concession privée~: le monopole d'État du télégraphe s'étendit au téléphone par décision administrative. En Angleterre, fidèle à sa doctrine de libéralisme économique, le gouvernement accorda des concessions à des compagnies privées, chacune opérant un réseau urbain dans sa zone d'activité. Mais le \emph{Post Office} perçut rapidement le danger que le téléphone représentait pour les recettes du télégraphe, et se réserva d'abord les lignes interurbaines (1899), puis nationalisa la totalité des réseaux en~1912 \parencite{bertho_histoire_1984}.

Le cas français est le plus instructif, parce qu'il met en jeu un mécanisme institutionnel absent des cas américain et prussien. En~1879, le gouvernement de la IIIe~République accorda des concessions téléphoniques à des sociétés privées. Trois candidats fusionnèrent avant la fin de l'année pour former la \emph{Société générale des téléphones} (SGT), dotée d'un capital de 25~millions de francs, qui prit en charge les réseaux de Paris, Lyon, Marseille et de quelques villes de province. Les termes de la concession prévoyaient un partage des responsabilités entre la SGT (les abonnés et les centraux) et l'administration (les câbles), dont la cohabitation fut difficile dès le début. Lorsque, en~1884, l'incertitude sur le renouvellement de la concession devint palpable, la SGT cessa d'investir~: le service se dégrada, les listes d'attente s'allongèrent, les abonnés d'affaires se plaignirent des tarifs \parencite{bertho_histoire_1984}. La tentative de la SGT de négocier secrètement avec le ministre des Finances un accord de \og compagnie fermière\fg{} pour trente ans provoqua un scandale qui précipita la nationalisation.

\dimmark{D7}{+}%
Ce qui est remarquable dans l'épisode français n'est pas la nationalisation elle-même~--- on l'a vue aussi en Angleterre~--- mais la justification qui la sous-tend. En~1837, quand le Parlement français avait débattu du statut du télégraphe, l'argument dominant était celui de la sécurité de l'État~: le télégraphe devait être un monopole public parce que la transmission rapide d'informations était une affaire de souveraineté (§\,\ref{sec:telegraph}). En~1889, quand la nationalisation du téléphone fut décidée, l'argument avait changé de nature. La chambre de commerce de Valenciennes, dans une délibération de~1887, formula la position qui allait prévaloir~: \og L'État ne doit pas se dessaisir de l'organisation de grands services d'utilité publique qui n'ont aucunement le caractère d'une industrie particulière, mais dont le but est de mettre à la portée de tous les moyens d'action capables de faciliter les progrès du travail et de la production nationale\fg{}\footnote{Cité par \textcite{bertho_histoire_1984}, p.~62. La délibération de la chambre de commerce de Valenciennes se prononçait \og contre toute concession de ligne téléphonique à une compagnie privée\fg{}.}. Le monopole de sécurité du télégraphe avait cédé la place à un monopole de \emph{service public}~: même forme institutionnelle, fondement politique entièrement différent. Dans les deux cas, l'État prenait le contrôle du réseau~; dans le premier, il le faisait pour garder le pouvoir~; dans le second, pour fournir un service universel.

La nationalisation française eut cependant un effet que ses promoteurs n'avaient pas anticipé. Le financement des réseaux téléphoniques fut organisé par le système des \og avances remboursables\fg{}~: les collectivités locales avançaient les fonds, remboursés sur les recettes du réseau. Ce mode de financement calqua la géographie du réseau sur la géographie administrative plutôt que sur la rationalité économique~: chaque canton, parfois chaque commune, eut son propre réseau, mal relié aux autres. En~1921, la France comptait 14\,000~centraux téléphoniques pour 7\,000~en Allemagne, alors que la population allemande était plus nombreuse et plus urbanisée \parencite{bertho_histoire_1984}. Le défaut principal du réseau~--- l'insuffisance des circuits interurbains~--- découlait directement du mode de financement. À partir de~1910, techniciens et grand public convergèrent dans le même diagnostic~: le téléphone français était en crise. Ce que Bertho nomme des \og vices de structure\fg{} ne relevaient pas de la technique mais de l'architecture institutionnelle. Des travaux ultérieurs ont confirmé que cette fragmentation avait des conséquences macroéconomiques mesurables~: \textcite{roller_telecommunications_2001}, étudiant 21~pays de l'OCDE entre~1970 et~1990 avec un modèle simultané qui traite la causalité inverse, ont montré que l'effet de l'infrastructure téléphonique sur la croissance n'apparaît qu'au-delà d'un seuil de pénétration d'environ 40\,\%~--- un seuil que la fragmentation française, avec ses 14\,000~centraux mal reliés, rendait structurellement difficile à atteindre.

\dimmark{D8}{+}%
La période de monopole du brevet Bell (1879--1894) avait été, pour les actionnaires, la plus profitable de l'histoire de l'industrie~: \textcite{gabel_early_1969}, exploitant les données du Walker Report de la FCC (1938), a calculé un rendement de près de quarante-six pour cent sur la période, avec des dividendes moyens de quinze pour cent. Mais cette rentabilité avait pour contrepartie la stagnation~: des tarifs élevés, des installations insuffisantes, un service limité au centre-ville des grandes villes, et un développement résidentiel, suburbain et rural à peine esquissé. Aux États-Unis, la fin de la protection par brevet, en~1894, ouvrit une période de concurrence violente qui produisit le taux de croissance le plus rapide de l'histoire de l'industrie. En~1900, plus de six mille compagnies indépendantes opéraient sur le territoire américain, et le parc téléphonique avait été multiplié par quatre en cinq ans \parencite{mueller_universal_1993}. Soixante pour cent des villes américaines de plus de cinq mille habitants avaient deux réseaux concurrents en~1904, ce qui forçait les abonnés au \og \emph{dual service}\fg{}~: souscrire aux deux pour pouvoir joindre tout le monde. Les indépendants couvraient les zones rurales et les petites villes que le Bell System avait négligées par calcul~: tant que les brevets protégeaient le monopole, il n'y avait pas de raison de desservir les marchés peu rentables\footnote{\textcite{mueller_universal_1993} a montré que la notion de \og \emph{universal service}\fg{}, telle que Vail l'employa pour la première fois dans le rapport annuel d'AT\&T de~1907, ne signifiait pas un téléphone dans chaque foyer mais l'interconnexion des réseaux en un système unifié. C'est la concurrence, non le monopole, qui avait étendu le réseau géographiquement~; le slogan de Vail visait à légitimer la reconsolidation, pas à décrire une réalité.}. En~1907, la situation était critique~: les profits de Bell baissaient. Un syndicat de banquiers, présidé par J.~P.~Morgan, prit le contrôle de la compagnie et rappela Vail.

Le second Vail réorganisa l'entreprise autour de deux piliers. Le premier était la recherche fondamentale~: il assigna aux laboratoires d'ingénierie (qui deviendraient les Bell Labs en~1925) la mission de construire une ligne transcontinentale entre New York et San Francisco, soit six mille kilomètres, alors que la ligne la plus longue (Boston-Chicago, 1893) ne dépassait pas mille neuf cents kilomètres. La solution vint de la triode de Lee de Forest (1906), dont AT\&T acheta les brevets en~1914 pour neuf cent mille dollars~: l'amplification du signal par le tube à vide rendit la transmission longue distance possible, et la ligne transcontinentale fut ouverte au trafic commercial le 15~janvier~1915 \parencite{bertho_histoire_1984}. Le second pilier était la régulation négociée. En~1913, face à une enquête antitrust, Vail proposa ce qui est resté dans l'histoire sous le nom de \emph{Kingsbury Commitment}~: AT\&T revendrait ses participations dans Western Union, cesserait d'acquérir des compagnies indépendantes sans autorisation, et ouvrirait ses lignes longue distance aux concurrents, moyennant un droit de passage. L'accord semblait pro-concurrentiel~; dans les faits, il institutionnalisa le monopole. La clause qui limitait les acquisitions fut interprétée de manière permissive~: AT\&T pouvait acheter un opérateur à condition d'en revendre un autre de taille équivalente, ce qui revenait à un échange de territoires plutôt qu'à une ouverture du marché. En~1924, AT\&T avait absorbé 223~des 234~compagnies indépendantes \parencite{weiman_preying_1994}. Au milieu des années~1920, l'ère de la concurrence était terminée.

\dimmark{D2}{*}%
Quatre pays, quatre trajectoires institutionnelles~--- le monopole privé intégré (États-Unis), le monopole d'État immédiat (Prusse), la concession suivie de nationalisation (France, Angleterre)~--- et un même résultat~: un seul opérateur par territoire. Le pattern observé dans les sept trajectoires du télégraphe (§\,\ref{sec:telegraph}) se confirme, ce qui renforce l'hypothèse que la convergence vers le monopole est un fait structurel des réseaux de communication, pas un accident institutionnel. Mais le mécanisme est chaque fois différent~: rendements de réseau pour le télégraphe, intégration verticale pour le Bell System, décision administrative pour la Prusse, défaillance de la concession pour la France. Ce qui varie n'est pas le résultat mais le chemin~--- et le chemin, à son tour, détermine les pathologies du monopole résultant (fragmentation française, prédation américaine, rigidité prussienne).


\subsection{La commutation~: du calcul dans le réseau}\label{subsec:commutation}

Le téléphone, à la différence du télégraphe, exigeait une opération que le réseau télégraphique ne connaissait pas~: la commutation. Un message télégraphique était adressé à une station, où un opérateur le transcrivait et le remettait au destinataire~; le routage se faisait par la géographie des lignes, pas par un geste de connexion en temps réel. Un appel téléphonique supposait au contraire qu'un circuit physique fût établi entre deux abonnés, pour la durée de la conversation, parmi des milliers ou des millions de connexions possibles. Ce circuit devait être créé à la demande, maintenu le temps de l'échange, puis défait. La personne qui accomplissait ce geste~--- chercher le numéro demandé sur le tableau, enfoncer une fiche dans la prise correspondante, surveiller la fin de l'appel~--- était l'opératrice de commutation.

\dimmark{D4/D7}{+}%
Le geste paraissait simple~; sa complexité croissait quadratiquement avec le nombre d'abonnés. Pour $N$~abonnés, le nombre de connexions possibles est $N(N-1)/2$~: un réseau de mille abonnés suppose près d'un demi-million de circuits potentiels. À mesure que le nombre d'abonnés augmentait, les tableaux de commutation s'étendaient, le nombre d'opératrices croissait, et la qualité du service se dégradait, parce que chaque appel devait transiter par un nombre croissant d'intermédiaires. Le \emph{switching} n'était pas une tâche de transmission~; c'était une tâche de traitement de l'information~--- du routage, du calcul implicite~--- intégrée au c\oe{}ur du réseau de transmission. C'est la première fois dans le parcours du chapitre que les fonctions de TRANSMIT et de COMPUTE apparaissent inséparables dans un même système opérationnel.

L'automatisation du \emph{switching} fut tentée dès~1891, lorsque Almon Brown Strowger, un entrepreneur de pompes funèbres de Kansas City qui soupçonnait une opératrice de détourner ses appels vers un concurrent, breveta un commutateur mécanique à pas de progression (brevet US n\textsuperscript{o}~447,918). Le système de Strowger permettait à l'abonné de composer lui-même le numéro à l'aide d'un cadran rotatif, chaque impulsion actionnant un sélecteur électromécanique au central. La technologie fonctionnait~; quelques compagnies indépendantes l'adoptèrent dans les années~1900. Mais AT\&T, qui possédait le plus grand réseau et employait le plus grand nombre d'opératrices, n'adopta le commutateur mécanique qu'à partir de~1919, et la diffusion prit soixante ans supplémentaires pour couvrir l'ensemble du réseau.

\dimmark{D7/D10}{+}%
Ce délai est un fait surprenant au regard du protocole de ce chapitre~: une technologie disponible depuis les années~1890, manifestement capable de réduire les coûts de main-d'\oe{}uvre, et déployée par le concurrent le mieux équipé pour l'adopter, mit près d'un siècle à se généraliser. \textcite{feigenbaum_organizational_2024} ont montré que l'explication ne résidait ni dans un défaut de la technologie ni dans une résistance syndicale, mais dans les interdépendances organisationnelles entre la commutation et le reste du système Bell. Les opératrices n'étaient pas seulement des connectrices de circuits~; elles étaient le pivot d'un système de production complexe qui s'était construit autour d'elles. Elles comptaient les appels pour la facturation~; elles géraient les plaintes en temps réel~; elles fournissaient les renseignements et les services d'urgence~; elles connaissaient les numéros des abonnés réguliers. Automatiser le \emph{switching} exigeait de réinventer chacune de ces fonctions adjacentes~: créer des compteurs automatiques, des services de renseignement séparés, des procédures d'urgence sans opérateur, des annuaires imprimés plus complets. La commutation était ce que la théorie de l'architecture organisationnelle appelle une \emph{tâche intégrale} \parencite{ulrich_role_1995}~--- une tâche qui interagit avec toutes les autres et dont la modification implique la reconfiguration de l'ensemble.

\emergence{D7/D10~--- L'automatisation d'une tâche apparemment simple (la commutation) prend soixante ans parce que cette tâche est \emph{intégrale} au système de production. Le \emph{switching} portait de l'intelligence organisationnelle, pas seulement de la connexion physique. Hypothèse abductive~: la résistance à l'automatisation est proportionnelle au degré d'intégralité de la tâche dans le système. \questionch{2}{D7~--- Le concept de tâche intégrale s'applique-t-il aux vagues d'automatisation ultérieures (computing, robotique, IA)~?}}


\subsection{Les demoiselles du téléphone}\label{subsec:demoiselles}

\dimmark{D4}{*}%
Les premières opératrices de commutation furent, à New York en~1879, de jeunes garçons recrutés parmi les porteurs de dépêches télégraphiques. Ils se révélèrent, selon les termes de l'époque, \og trop remuants\fg{} pour se plier à la discipline physique du tableau de commutation~: rester assise, bridée par le casque d'écoute et le microphone, incapable de se lever ni de se rasseoir, tenue à une courtoisie sans faille envers des abonnés souvent irrités. Dès~1880, la \emph{Société générale des téléphones} à Paris n'embauchait plus que des femmes~; AT\&T fit de même aux États-Unis. En~1920, les opératrices de téléphone constituaient l'un des emplois féminins les plus répandus en Amérique~: environ 4\,\% des jeunes femmes blanches nées aux États-Unis qui travaillaient étaient opératrices, et AT\&T était le plus grand employeur de femmes du pays, avec des opératrices représentant plus de la moitié de sa main-d'\oe{}uvre \parencite{feigenbaum_answering_2024}.

Le profil de recrutement était précis~: diplômée du \emph{high school}, au moins dix-huit ans, vivant chez ses parents ou un proche, en bonne santé physique, dotée d'une voix agréable, d'une dextérité manuelle, d'une capacité à calculer rapidement et d'un tempérament stable, \og \emph{not easily ruffled by irritable customers}\fg{} \parencite{erickson_telephone_1946}. Les opératrices étaient presque exclusivement blanches, nées aux États-Unis~; AT\&T ne recruta pas de femmes noires avant~1940 \parencite{green_race_2001}. Le poste était mieux payé que le travail en usine, offrait une interaction humaine que le travail de bureau ne garantissait pas, mais imposait un rythme intense et un taux de rotation de l'ordre de 40\,\% par an \parencite{feigenbaum_answering_2024}.

Le parallèle avec le cas du calcul mécanisé (§\,\ref{subsec:zoom_tabulating}) est frappant. Au \emph{Nautical Almanac}, des calculateurs qualifiés, autonomes, travaillant à domicile, avaient été remplacés par des opératrices moins payées, recrutées dans les catégories inférieures de la fonction publique, supervisées de près. Au central téléphonique, le même schéma se reproduit~: une tâche exigeant des qualités \og spécifiquement féminines d'habileté, de patience, de résistance nerveuse\fg{} \parencite{bertho_histoire_1984} est confiée à des jeunes femmes dont la principale vertu, du point de vue du recruteur, est d'accepter un salaire inférieur à celui que des hommes auraient exigé pour le même travail. La féminisation et la compression salariale documentées par \textcite{davies_womans_1982} et \textcite{strom_beyond_1992} pour les machines de bureau se reproduisent dans le téléphone, ce qui renforce l'hypothèse que ce pattern est structurel et non contingent.

\dimmark{D4}{*}%
L'automatisation de la commutation, entre~1920 et~1940, élimina la majorité de ces emplois, une ville à la fois. \textcite{feigenbaum_answering_2024}, exploitant la variation dans le calendrier d'adoption du commutateur mécanique à travers près de trois mille villes américaines, ont montré que le nombre de jeunes opératrices chutait de 50 à 80\,\% dans les années qui suivaient l'installation du système automatique. L'effet sur les opératrices en poste était sévère~: une décennie après l'automatisation, celles qui travaillaient encore occupaient des emplois moins bien payés que leurs homologues dans les villes non automatisées, et les plus âgées d'entre elles avaient quitté le marché du travail. L'automatisation détruisait non seulement des emplois mais des carrières~: le \emph{switching} était l'un des rares métiers féminins de l'époque qui offrait une perspective de long terme, et sa disparition fermait cette possibilité.

\dimmark{D4}{+}%
Mais le résultat le plus significatif de l'étude concerne les cohortes suivantes. Les jeunes femmes qui entraient sur le marché du travail après l'automatisation ne connaissaient pas de baisse de leur taux d'emploi global. Les emplois perdus dans la commutation étaient compensés par une croissance de l'emploi dans le secrétariat (emplois de qualification comparable) et dans la restauration (emplois de qualification inférieure). Plus remarquable encore, la croissance de l'emploi secrétarial se concentrait dans des industries qui n'avaient pas employé de secrétaires auparavant~--- de nouveaux pairages occupation-industrie, ce que \textcite{autor_new_2024} appellent du \og \emph{new work}\fg{}. L'automatisation n'avait pas seulement déplacé les travailleuses d'un secteur vers un autre~; elle avait engendré, indirectement, la création d'usages nouveaux du travail féminin dans des secteurs qui n'en avaient pas eu besoin jusque-là. Ce mécanisme de réinstauration des tâches, théorisé par \textcite{acemoglu_race_2018}, était cependant conditionné~: dans les villes manufacturières, où les emplois féminins en col blanc étaient rares, et dans les villes frappées le plus durement par la Grande Dépression, l'automatisation réduisait l'emploi sans le compenser. La demande agrégée est une condition nécessaire pour que le mécanisme de réinstauration opère.

\emergence{D4~--- L'automatisation de la commutation (1920--1940) confirme et étend le pattern du \emph{Nautical Almanac}~: la mécanisation réorganise le travail plutôt qu'elle ne le supprime. Les incumbentes souffrent~; les cohortes suivantes se réallouent. Le mécanisme de compensation passe par la création de \og \emph{new work}\fg{} (Autor et al.~2023) dans des industries qui n'employaient pas de secrétaires auparavant. Mais ce mécanisme est conditionné par la demande agrégée~: en période de contraction, l'automatisation réduit l'emploi sans le compenser. \questionch{2}{D4~--- La conditionnalité du mécanisme de réinstauration à la demande agrégée a-t-elle une implication pour les vagues d'automatisation contemporaines (IA)~?}}


\subsection{La bobine Pupin et la matérialité du réseau}\label{subsec:pupin}

\dimmark{D3/0-TER}{*}%
Un dernier fait mérite d'être relevé avant de clore le cas, parce qu'il concerne la dimension matérielle que le récit du monopole et de l'automatisation a laissée en arrière-plan. En~1900, la communication longue distance par téléphone se heurtait à un obstacle physique~: l'affaiblissement du signal sur les lignes de cuivre. La ligne Boston-New York, inaugurée en~1884, représentait à peu près la limite pratique. Un quart du capital investi par les compagnies téléphoniques allait au fil de cuivre, et l'allongement des distances exigeait des conducteurs de section croissante, donc de coût croissant. Michael Pupin, professeur de physique à Columbia, breveta en~1900 une bobine d'inductance qui, insérée à intervalles réguliers sur la ligne, réduisait l'affaiblissement et permettait de doubler la distance de communication à section de cuivre égale, ou de réduire de moitié la section nécessaire pour une distance donnée. AT\&T acheta les brevets de Pupin pour un total de 455\,000~dollars entre~1901 et~1917 \parencite{brittain_turning_1970}. L'économie de cuivre réalisée a été estimée à 100~millions de dollars pour le premier quart du vingtième siècle \parencite{brittain_turning_1970}, soit un ratio de 1~pour~200 entre le coût du brevet et l'économie de matière qu'il permit.

Ce ratio illustre un mécanisme de substitution du capital intellectuel (le brevet, la connaissance mathématique de la propagation des ondes) au capital physique (le cuivre). C'est un fait de même nature que la substitution K/L documentée dans le cas du télégraphe (§\,\ref{sec:telegraph}), mais il opère sur un autre facteur~: non pas le travail mais la matière première. La bobine Pupin ne transmettait rien~; elle permettait de transmettre avec moins de cuivre. Le fait que le poste de dépense dominant du réseau téléphonique fût le cuivre, pas l'énergie ni le travail, signale que la matérialité du réseau de transmission pesait lourd dans l'économie du système, même si cette matérialité restait invisible pour l'abonné.

\dimmark{0-TER}{+}%
Quant à la dépendance énergétique, le cas du téléphone marque une transition qualitative par rapport au télégraphe, même si elle est discrète. Le télégraphe de Morse tirait son énergie de piles voltaïques, remplacées localement quand elles s'épuisaient~; la consommation était intermittente et décentralisée. Le téléphone, dans ses premières années, fonctionnait sur le même modèle~: chaque poste avait sa propre pile. Mais l'introduction de la \og batterie centrale\fg{}, où une pile humide unique, placée au central, alimentait l'ensemble des postes connectés, transforma le réseau téléphonique en un système à consommation permanente et centralisée\footnote{L'invention de la batterie centrale est attribuée à Pavel Golubitsky en Russie (années~1880). Elle fut rapidement adoptée par l'ensemble des réseaux urbains parce qu'elle simplifiait la maintenance et permettait l'extension à grande échelle.}. C'est un basculement modeste en valeur absolue~--- la consommation d'un central téléphonique de~1900 était négligeable en regard de celle d'une centrale électrique~--- mais c'est le premier pas dans la séquence d'escalade qui mènera, par le tube à vide (§\,\ref{sec:wireless}), le mainframe (§\,\ref{sec:computing}) et le datacenter (§\,\ref{sec:contemporain}), à une dépendance énergétique d'un tout autre ordre de grandeur.


\subsection*{Bilan~: ce que le téléphone ajoute à la traversée}\label{subsec:bilan_telephone}

Le cas du téléphone confirme certains patterns établis dans les cas précédents et en casse d'autres. La convergence vers le monopole se confirme~: quatre pays, quatre trajectoires institutionnelles, un même résultat. Mais le mécanisme est nouveau (intégration verticale et contrôle de l'interconnexion), ce qui renforce l'hypothèse que c'est le mode de contrôle du système, pas la technologie, qui détermine la structure de marché. La féminisation et la compression salariale se confirment (les demoiselles du téléphone prolongent les calculatrices du \emph{Nautical Almanac}), mais l'automatisation ajoute un résultat que le cas du calcul mécanisé n'avait pas produit~: les cohortes suivantes se réallouent, et la demande de travail se reconstitue par la création de tâches nouvelles dans des industries qui n'en avaient pas eu l'usage.

Ce qui est propre au téléphone, et qui ne pouvait pas émerger des cas précédents, tient en quatre observations. La première est la séquence temporelle du déploiement~: le téléphone commence comme commodité urbaine locale (\textcite{ploeckl_market_2023}), puis la concurrence post-brevets (1894--1907) étend le réseau géographiquement (\textcite{gabel_early_1969}), et c'est seulement après~1907 que la reconsolidation sous Vail en fait un outil de coordination organisationnelle à longue distance. Le déplacement du canal d'action~--- de la communication locale vers la coordination intra-organisationnelle~--- est un processus de trente ans, pas un fait contemporain de l'invention\footnote{\textcite{norton_transaction_1992}, utilisant le stock initial de téléphones en~1957 comme prédicteur de la croissance~1957--1977, a formalisé le canal par lequel le téléphone agit sur l'économie~: la réduction des coûts de transaction (au sens de \textcite{leff_externalities_1984}). \textcite{portes_determinants_2005} utilisent le trafic téléphonique bilatéral entre pays comme proxy de friction informationnelle dans un modèle de gravité des flux de capitaux~--- c'est l'un des rares usages économétriques du téléphone comme \emph{variable} d'information. \textcite{bloom_distinct_2014} ont montré que les technologies de communication centralisent les décisions dans la firme tandis que les technologies d'information les décentralisent~; cette distinction, qui recoupe celle que le présent chapitre construit entre le télégraphe et le téléphone, sera reprise au chapitre~2.}.

La deuxième observation est la communication domestique. Le téléphone est la première technique de l'information du parcours de ce chapitre dont l'utilisateur final n'est pas une entreprise, un État ou un militaire, mais un particulier chez lui. \textcite{hardy_role_1980}, dans une analyse cross-country couvrant 60~nations entre~1960 et~1973, a trouvé que les téléphones résidentiels contribuaient davantage à la croissance économique que les téléphones d'affaires (coefficient de~0,249 contre~0,097 avec la consommation d'énergie comme indicateur de développement)~--- un résultat qu'il qualifie lui-même de \og \emph{quite unexpected}\fg{}. Ce résultat s'éclaire sous Ploeckl~: si le téléphone est d'abord une commodité locale, c'est l'usage résidentiel qui capte l'essentiel de la valeur initiale. \textcite{martin_culture_1998} a montré que les femmes accédèrent au téléphone par l'intermédiaire des hommes d'affaires qui faisaient connecter leur domicile, et développèrent des usages~--- la sociabilité, la coordination domestique~--- que l'industrie considérait comme parasitaires avant de s'y adapter. La diffusion elle-même n'est pas monotone~: \textcite{fischer_technologys_1987} a documenté la perte d'un million de téléphones dans les fermes américaines entre~1920 et~1940, une dévolution technologique qui montre que l'adoption d'une technique de l'information peut reculer quand les conditions économiques se dégradent.

La troisième est la nature intégrale de la commutation~: le \emph{switching} n'était pas une tâche isolable mais le pivot d'un système de production complexe, et sa résistance à l'automatisation pendant soixante ans en est la preuve. La quatrième est la transition du monopole de sécurité au monopole de service public~: en~1837, l'État prenait le contrôle du réseau pour garder le pouvoir~; en~1889, pour fournir un service universel.

\emergence{D10~--- Le cas de Gray révèle un mécanisme plus précis que le \emph{focusing device} de \textcite{rosenberg_perspectives_1976}. Ce n'est pas seulement que le paradigme en place dirige l'attention vers certains problèmes~; c'est que le système métrique du paradigme~--- le nombre de messages par fil, la vitesse, le coût par mot~--- rend impossible de \emph{mesurer la valeur} d'innovations qui opèrent dans un autre registre. Gray ne manque pas d'intelligence~; il manque d'unités~--- et cette absence d'unités n'est pas une carence cognitive individuelle mais le résultat d'un cadre institutionnel (l'industrie télégraphique, son modèle d'affaires fondé sur la facturation au mot, ses contrats d'emploi qui rémunèrent l'optimisation du multiplexage) qui n'a pas besoin de telles unités pour fonctionner et qui en décourage activement la production. Que Bell, lui-même, ait dû lutter contre ses commanditaires Hubbard et Sanders pour se détourner du multiplexage le confirme indirectement~: ce ne sont pas seulement les esprits qui sont structurés par le paradigme, ce sont aussi les contrats qui les financent. Le téléphone ne produit pas de \og messages\fg{} au sens du télégraphe, et sa valeur~--- la communication de proximité, la sociabilité, le non-codifiable~--- n'existe pas dans la métrique du paradigme en place. \questionch{2}{D10~--- Ce mécanisme d'occultation de la valeur par le système métrique du paradigme se reproduit-il dans les cas ultérieurs (téléphone $\rightarrow$ internet, PC $\rightarrow$ smartphone)~?}}

Le tableau~\ref{tab:telephone-dimensions} récapitule les observations dimensionnelles du cas.

\begin{table}[p]
\caption{Le téléphone au prisme des dix dimensions~: observations, niveaux de certitude et questions ouvertes}\label{tab:telephone-dimensions}
\centering\footnotesize
\renewcommand{\arraystretch}{1.15}
\begin{tabularx}{\textwidth}{@{}cXcp{4.5cm}@{}}
\toprule
\textbf{Dim.} & \textbf{Observations du cas §\,\ref{sec:telephone}} & & \textbf{Question pour la suite} \\
\midrule
D1 & Le téléphone transmet un signal continu (la voix), non un signal discret (le code Morse). L'information transmise n'est pas réductible à un alphabet~; sa valeur réside dans ce que le codage détruirait (intonation, hésitation, crédibilité). & + & Cette distinction codifiable/non-codifiable détermine-t-elle le canal par lequel l'ICT affecte l'économie~? \\[3pt]
D2 & Convergence vers le monopole par un troisième mécanisme (intégration verticale + contrôle de l'interconnexion). Confirmé dans quatre contextes institutionnels distincts (US, FR, UK, DE). Prédation documentée par \textcite{weiman_preying_1994}. & * & Le mécanisme de monopole dépend-il de la technologie, de l'institution, ou du hasard historique~? \\[3pt]
D3 & Ratio Pupin~: \$455k de brevets pour \$100M de cuivre économisé (1:200). Substitution du capital intellectuel au capital physique (matière, pas travail). & * & Ce ratio de substitution se maintient-il avec les technologies ultérieures~? \\[3pt]
D4 & Féminisation et compression salariale~: même pattern que le calcul mécanisé. L'automatisation (1920--1940) élimine les emplois d'opératrices mais ne réduit pas l'emploi total des cohortes suivantes \parencite{feigenbaum_answering_2024}. Réallocation vers le secrétariat et la restauration. \og \emph{New work}\fg{} dans des industries nouvelles. & * & Le mécanisme de réinstauration est-il conditionné par la demande agrégée dans tous les cas~? \\[3pt]
D5 & 48\,000 abonnés (1880) $\rightarrow$ 801\,000 (1900) $\rightarrow$ 7~millions (1906). Transition tarifaire~: du tarif au message (comme le télégraphe) au forfait. Les téléphones résidentiels contribuent davantage à la croissance que les business phones (\textcite{hardy_role_1980}~: coeff.~0,249 vs 0,097). La diffusion n'est pas monotone~: un million de fermes perdent le téléphone entre~1920 et~1940 \parencite{fischer_technologys_1987}. & + & Le passage au forfait et la domestication modifient-ils la structure de la demande de communication~? \\[3pt]
D6 & Pas d'étude d'identification causale de l'effet du téléphone sur l'intégration des marchés. \textcite{ploeckl_market_2023} montre que le téléphone était d'abord une commodité urbaine locale, pas un outil d'intégration de marchés distants. L'absence reflète un déplacement progressif du canal d'action (local $\rightarrow$ longue distance $\rightarrow$ organisationnel). & * & Ce déplacement temporel du canal d'action se retrouve-t-il dans d'autres technologies de communication~? \\[3pt]
D7 & Le \emph{switching} est une tâche intégrale (Ulrich 1995)~: il interagit avec toutes les fonctions du système Bell. 60~ans pour automatiser. Transition justification du monopole~: sécurité (1837) $\rightarrow$ service public (1889). & + & Le concept de tâche intégrale s'applique-t-il aux vagues d'automatisation ultérieures~? \\[3pt]
D8 & Cycle~: monopole de brevet (1876--1894, stagnation, rendement 46\,\% pour les actionnaires, \textcite{gabel_early_1969}) $\rightarrow$ concurrence (1894--1907, 6\,000 compagnies, croissance la plus rapide de l'industrie) $\rightarrow$ reconsolidation (Morgan/Vail, 1907--1925). Le \emph{Kingsbury Commitment} (1913) institutionnalise le monopole sous couvert de régulation. Masse critique~: l'effet macro n'apparaît qu'au-delà de 40\,\% de pénétration \parencite{roller_telecommunications_2001}. & * & Ce cycle ouverture-concurrence-reconsolidation est-il intrinsèque aux réseaux de communication~? \\[3pt]
D9 & Le cuivre est le poste de dépense dominant du réseau (25\,\% du capital investi). La bobine Pupin réduit la dépendance au cuivre par substitution intellectuelle. & + & La géographie de la matière première (cuivre) contraint-elle la géographie du réseau~? \\[3pt]
D10 & Le système métrique du paradigme en place (messages/fil/temps) rend impossible de mesurer la valeur d'innovations d'un autre registre. Gray ne manque pas d'intelligence~; il manque d'unités. Le téléphone, contrairement au télégraphe, naît hors du réseau ferroviaire \parencite{ploeckl_market_2023}~: la séparation communication/transport est empiriquement confirmée. & * & L'occultation de la valeur par le système métrique du paradigme se reproduit-elle dans les vagues ultérieures~? \\[3pt]
\midrule
0-TER & Le cuivre domine la structure de coûts (matière). L'énergie de fonctionnement est faible mais en transition~: des piles individuelles (décentralisé, intermittent) à la batterie centrale (centralisé, permanent). Premier basculement dans la séquence d'électrification de TRANSMIT. & + & À quel moment le coût énergétique dépasse-t-il le coût matériel dans les réseaux de communication~? \\[3pt]
\bottomrule
\end{tabularx}
\renewcommand{\arraystretch}{1.0}
\end{table}



% =============================================================================